\documentclass[exercice,a4]{thedocument}%
\usepackage{texmaths-tikz}%
\usepackage{pstricks}%
\usepackage{pst-bezier}%
\usepackage{pstricks-add}%
\usepackage{pst-eucl}%
\usepackage{pst-func}%
\usepackage{pst-tree}%
\usepackage{pst-circ}%
\usepackage{tikz}%
\usepackage{pgfplots}%
\usepackage{tkz-tab}%
\usepackage{tikzpagenodes}%
\usepackage{tabularray}%
\usepackage{tabularx}%
\usepackage{multirow}%
\usepackage{multicol}%
\usepackage{scratch3}%
\usepackage{enumitem}%
\usepackage[np]{numprint}%
\usepackage{pas-tableur}%
\usepackage{listings}%
\usepackage{graphicx}%
\usepackage{adjustbox}%
\usepackage{needspace}%
\usepackage{xcolor}%
\usepackage{eurosym}%
\usepackage{diagbox}%
\usepackage{colortbl}%
\usepackage{multido}%
\usepackage{icomma}%
\usepackage[french]{babel}%
\startdatakeys
    \source{26GENMATG11}
    \BO{2026}
    \cycle{4}
    \competencesDuSocle{CA,RA,CO}
    \objectifsApprentissage{C4A1aa,C4A2ad,C4B2cb,C4D2aa,C4D3ac}
    \motsCles{programme de calcul,Scratch,développement,équation}
\enddatakeys
\definecolor{gris}{gray}{0.5}%
\providecommand{\wideparen}[1]{\overset{\frown}{#1}}%
\begin{document}%
\providecommand{\degres}{^\circ}%
Voici deux programmes de calcul :

\begin{center}
	\begin{minipage}{0.48\textwidth}
		\centering \textbf{Programme A}

		\begin{center}
			%Début Tikz
			\usetikzlibrary{positioning}
			\begin{tikzpicture}[every node/.style={draw,line width=1pt, fill=white, align=center, rectangle, inner xsep=16pt, inner ysep=6pt, minimum width=2.8cm,minimum height=8mm},> = stealth]
				\draw[->,line width=1pt] (-0.8,-0.4) -- (-2.5,-0.9);
				\draw[->,line width=1pt] (0.8,-0.4) -- (2.5,-0.9);
				\draw[->,line width=1pt] (-2.8,-1.8) -- (-0.8,-2.3);
				\draw[->,line width=1pt] ( 2.8,-1.8) -- ( 0.8,-2.3);
				\node at(0,0) {Choisir un nombre};
				\node at (-2.8,-1.4) {Ajouter 6};
				\node at (2.8,-1.4) {Soustraire 4};
				\node[below] at (0,-2.4) {Multiplier les\\deux nombres obtenus};

			\end{tikzpicture}
			%Fin Tikz
%			%Début PSTricks
%			\psset{unit=1cm,arrowsize=2pt 3}
%			\begin{pspicture}(-3.5,0)(3.5,4)
%				\rput(0,3.5){\fbox{Choisir un nombre}}\psline{->}(-0.6,3.2)(-2.1,2.5)
%				\rput(-2,2.2){\fbox{Ajouter 6}}
%				\rput(2,2.2){\fbox{Soustraire 4}}\psline{->}(0.8,3.2)(2,2.5)
%				\rput(0,0.9){Multiplier les}\psline{->}(-2,1.9)(-0.8,1.1)
%				\rput(0,0.6){deux nombres obtenus}\psline{->}(2,1.9)(0.8,1.1)
%				\psframe(-2,0.4)(2,1.1)
%			\end{pspicture}
%			%Fin PSTricks
		\end{center}
	\end{minipage}
	\hfill
	\begin{minipage}{0.48\textwidth}
		\centering \textbf{Programme B}\\[0.3em]
		\begin{flushleft}
			\textbullet~\ Choisir un nombre.\\
			\textbullet~\ Calculer le carré de ce nombre.\\
			\textbullet~\ Soustraire 9.
		\end{flushleft}
	\end{minipage}
\end{center}

\begin{enumerate}
	\item On choisit 1 comme nombre de départ. Vérifier que le résultat obtenu avec le programme A est $- 21$.
	\item On choisit 10 comme nombre de départ. Calculer le résultat obtenu avec le programme B.
	\item Donner tous les nombres de départ possibles qui permettent d'obtenir 16 avec le programme B.
	\item \begin{minipage}[t]{0.35\linewidth}
		\textbf{(Question algorithmique)}

		Le programme ci-contre a été conçu avec le logiciel Scratch.

		Recopier et compléter sur la copie les lignes 3, 4 et 5 pour qu'il affiche le résultat obtenu avec le programme A lorsqu'un nombre de départ est saisi.

		\emph{Aucune justification n'est attendue.}
	\end{minipage}\hfill
		\begin{scratch}[scale=0.9,num blocks]
			\blockinit{Quand \greenflag est cliqué}
			\blocksensing{demander \ovalnum{choisir un nombre} et attendre}
			\blockvariable{mettre \selectmenu{Valeur 1} à \ovaloperator{\ovalvariable{réponse} + \ovalnum{}}}
			\blockvariable{mettre \selectmenu{Valeur 2} à \ovaloperator{\ovalvariable{réponse} - \ovalnum{}}}
			\blockvariable{mettre \selectmenu{Résultat} à \ovalnum{} * \ovalnum{}}
			\blocklook{dire \ovaloperator{regrouper\ovalnum{ Le résultat du programme A est} et {\ovalvariable{Résultat}}}}
		\end{scratch}
	%\begin{tabular}{|c|l|}
	%\hline
	%1 & Quand [drapeau vert] est cliqué \\
	%2 & demander \texttt{(Choisir un nombre)} et attendre \\
	%3 & mettre \texttt{Valeur 1} à \texttt{(réponse) + ( \ \ )} \\
	%4 & mettre \texttt{Valeur 2} à \texttt{(réponse) -- ( \ \ )} \\
	%5 & mettre \texttt{Résultat} à \texttt{( \ \ ) * ( \ \ )} \\
	%6 & dire \texttt{(regrouper(Le résultat du programme A est) et (Résultat))} \\
	%\hline
	%\end{tabular}
	%
	%\emph{(Lignes 3, 4 et 5 à compléter par le candidat avec les valeurs manquantes.)}

	\item On choisit $x$ comme nombre de départ. Montrer que le résultat obtenu avec le programme A est $x^2 + 2x - 24$.
	\item On cherche quel nombre de départ choisir pour que les programmes A et B donnent le même résultat. Écrire une équation permettant d'obtenir ce nombre de départ, puis la résoudre.
\end{enumerate}
\medskip
\showthesource
\end{document}
